%% AISec 2026 submission — ACM sigconf, double-blind review build.
%%
%% Compile on Overleaf:  Menu -> Compiler -> pdfLaTeX, then Recompile.
%% For the camera-ready, delete "anonymous,review" from the options below and
%% fill in \author / \affiliation, plus the real artifact URL in \S Availability.
\documentclass[sigconf,anonymous,review]{acmart}

%% AISec is not an archival-rights venue for the review copy; these three lines
%% suppress the "unpublished working draft" banner and the ACM reference block
%% that otherwise appear on a submission build.
\settopmatter{printacmref=false}
\renewcommand\footnotetextcopyrightpermission[1]{}
\pagestyle{plain}

\usepackage{booktabs}
\usepackage{balance}

\begin{document}

\title{Trust Boundaries and the Limits of Pattern-Based Prompt Injection
Detection}

%% Anonymous build: acmart suppresses these automatically under [anonymous].
\author{Anonymous Author(s)}
\affiliation{%
  \institution{Submission under double-blind review}
  \country{}
}

\begin{abstract}
We evaluate a two-tier prompt injection guardrail, an Aho--Corasick blocklist
followed by a fine-tuned DeBERTa-v3 classifier, and report a sequence of
measurements that each revise its apparent performance downward. On the
271-sample corpus its rules were tuned against, the system reaches 1.000 recall
at zero false positives. Held out from that same corpus it reaches 0.789. On
Open-Prompt-Injection it reaches 0.40, and on BIPIA 0.48. The shortfall is not a
tuning deficit. Attacks that announce themselves with override vocabulary are
caught almost perfectly, and attacks that inject a well-formed instruction
carrying no override phrasing are largely invisible. We then show that one bit of
deployment context, whether text arrived as a user turn or as data the
application is processing, moves recall on Open-Prompt-Injection from 0.40 to
1.00 at a 0.6\% false positive rate, because the same sentence is a legitimate
request on one side of that boundary and an attack on the other. Two further
results have practical consequences for anyone deploying such a system. INT8
dynamic quantization silently reduced the classifier tier to a constant
\textsc{safe} prediction, which a fail-open policy rendered invisible. And on
NotInject the deterministic tier blocks 0 of 339 benign trigger-word prompts
while the learned tier blocks 137, which makes the learned tier unshippable as a
blocking stage. Artifacts and every measurement script are public.
\end{abstract}

\begin{CCSXML}
<ccs2012>
<concept>
<concept_id>10002978.10003014</concept_id>
<concept_desc>Security and privacy~Network security</concept_desc>
<concept_significance>500</concept_significance>
</concept>
<concept>
<concept_id>10010147.10010257</concept_id>
<concept_desc>Computing methodologies~Machine learning</concept_desc>
<concept_significance>300</concept_significance>
</concept>
</ccs2012>
\end{CCSXML}
\ccsdesc[500]{Security and privacy~Network security}
\ccsdesc[300]{Computing methodologies~Machine learning}

\keywords{prompt injection; LLM security; guardrails; evaluation methodology;
trust boundaries; over-defense}

\maketitle

\section{Introduction}

Prompt injection is the top entry in the OWASP list for LLM applications, and the
defenses deployed against it are mostly detectors: a filter sits in front of the
model and decides whether a prompt is hostile. Published detectors converge on
two designs, regular expressions over attack vocabulary and fine-tuned
transformer classifiers~\cite{munirathinam2026beyond}. We built one of each,
wired them in series, and evaluated the result far past the point where the
numbers stopped flattering us.

This paper is a report on what that evaluation found. We do not present a new
detection algorithm. We present measurements of a conventional one, taken under
conditions that most published evaluations of this kind do not apply, and the
measurements disagree with each other in ways that turn out to be informative.

The headline is a sequence. Measured on the corpus its blocklist was tuned
against, our Tier~1 rule engine achieves perfect recall at zero false positives.
Held out from that corpus it achieves 0.789. On the standard external
benchmark~\cite{liu2024formalizing} it achieves 0.40. On an indirect-injection
benchmark~\cite{yi2023bipia} it achieves 0.48. Each number is honest; they differ
because they measure different things, and only the last two say anything about a
deployment.

What separates them is not attack difficulty. It is vocabulary. Our detector,
like the blocklists it is representative of, keys on the language of override:
\emph{ignore previous instructions}, \emph{disregard the above}, \emph{you are
now DAN}. An attacker who appends a well-formed instruction and no override
phrasing at all, which is exactly what three of the five strategies
in~\cite{liu2024formalizing} do, walks through. On those three strategies our
recall was zero.

The fix is not more vocabulary. It is context. In the threat model
of~\cite{liu2024formalizing} the injected instruction lands in the \emph{data} an
application processes, not in the user's own turn. Our detector was screening
both identically, which is incoherent: our benign corpus contains ``Summarize
this document for me,'' a legitimate user request that is indistinguishable, as a
string, from an injected one. Once callers declare which side of the trust
boundary a text arrived on, recall on~\cite{liu2024formalizing} moves from 0.40
to 1.00 at a 0.6\% false positive rate, with the user path unchanged.

We claim four contributions. First, a quantification of in-sample inflation for
this class of detector: 1.000 in-sample, 0.789 held out, 0.48 external, on the
same system. Second, evidence that trust-boundary conditioning is worth more than
vocabulary: one declared bit moves recall by 60 points where added keywords moved
it by 2.6. Third, a characterization of two opposing failure modes, where the
deterministic tier is precise but narrow (0.000 false positive rate on
NotInject~\cite{li2024injecguard}, 0.48 recall on BIPIA) and the learned tier is
neither precise nor broad (0.404 false positive rate, 0.13 recall). Fourth, a
negative result that generalizes past our system: pattern matching over
instruction surface form does not transfer across instruction phrasings.

\section{Background and related work}

Liu et al.~\cite{liu2024formalizing} give the formalization we use throughout,
along with the benchmark that generates injected prompts by splicing an injected
task's instruction and data into a target task's data under five strategies. Yi
et al.~\cite{yi2023bipia} introduce BIPIA for the indirect case, where the
payload is embedded in retrieved content rather than supplied directly. Both
release code, and we reproduce their construction rather than approximating it.

Li and Liu~\cite{li2024injecguard} identify over-defense as the failure mode that
actually breaks guard models in deployment, and release NotInject: 339 benign
prompts seeded with attack trigger words, on which state-of-the-art guards drop
to near-random accuracy. Their framing is why we report false positive rate on
adversarially benign text as a first-class number rather than as a footnote.

Nasr et al.~\cite{nasr2025attacker} evaluate twelve published defenses under
adaptive attack and exceed 90\% success on most, noting that the majority had
originally reported near-zero attack success. That paper is the reason our corpus
has a held-out split at all. Hackett et al.~\cite{hackett2025bypassing}
demonstrate character-injection evasion against production guardrails with up to
100\% success, which is a direct attack on the normalization our Tier~1 depends
on; we discuss this in Section~\ref{sec:limitations} as an evaluated-nowhere gap
rather than pretending coverage.

Closest to our system is Corll~\cite{corll2026mirror}, who argues that a first
screening layer should be fast, deterministic, non-promptable and auditable, and
reports 95.97\% recall on a genuine holdout with a sparse character $n$-gram SVM.
We agree with the design argument and adopt the holdout discipline.
Munirathinam~\cite{munirathinam2026beyond} names the regex-plus-classifier
pairing as the converged open-source architecture and predicts both failure modes
we go on to measure, so we make no novelty claim for the architecture itself.

\section{System}

The service exposes one endpoint and screens in two stages. Tier~1 normalizes the
input, then scans it with an Aho--Corasick automaton over 197 blocklist phrases
and 14 structural regular expressions. Normalization strips zero-width
characters, maps
mathematical, fullwidth, circled and small-capital Unicode to ASCII, applies
NFKD, removes combining marks, and decodes base64, hex and URL payloads when the
decoded text contains injection-relevant keywords.

Tier~1 scans twice. The first pass runs over the normalized text with separators
intact, so that literal markers such as \texttt{<|im\_start|>} and
\texttt{safety\_mode=off} still match. The second runs over a de-obfuscated
variant that maps leetspeak digits and symbols back to letters and strips
intra-word separators. The second pass is what catches \texttt{1gn0r3 prev10us
1nstruct10ns}, \texttt{!gnore prev!ous} and \texttt{instr-uction-s} without the
first pass losing its structural tokens. Median Tier~1 latency is 0.09~ms.

Tier~2 is \texttt{protectai/deberta-v3-base-prompt-injection-v2} exported to
ONNX, consulted only for prompts Tier~1 clears, with a block threshold of 0.75
and median latency 24~ms.

\subsection{Trust boundaries}

The endpoint takes an optional \texttt{source} field with two values. Under
\texttt{user}, the default, text is the end user's own turn and instructions in
it are legitimate. Under \texttt{data}, text is content the application is
processing, retrieved documents or tool output or third-party payloads, and an
instruction addressed to a model is itself the signal, because data is meant to
be read and not obeyed. Callers already possess this bit: an application passing
a retrieved email knows it is not passing a user turn.

Under \texttt{data} we additionally apply 12 patterns matching instruction shape
rather than attack vocabulary: forward references to following content,
output-format constraints such as \emph{answer with positive or negative}, role
assignment such as \emph{you are a tool}, sentence-initial task verbs, and
assistant-directed question forms.

\section{Evaluation}

We use an in-house corpus of 271 samples, 190 attacks across ten classes and 81
benign, plus three external benchmarks. All measurements run against the
containerized service on commodity hardware with CPU inference.

\subsection{In-sample versus held out}

Tier~1's blocklist was originally tuned the way such lists usually are: by
reading the attacks it missed on the evaluation corpus and adding terms drawn
from those samples. Recall measured afterwards on the same corpus is
memorization. To get a defensible number we split the corpus 60/40, stratified by
class and seeded, reconstructed the pre-tuning ruleset, and refit using
train-split misses only. Table~\ref{tab:heldout} shows the result.

\begin{table}[t]
\caption{Tier~1 recall by rule configuration. The third row is not a result; its
rules saw the test half, and it is shown only to quantify the optimism that
produces.}
\label{tab:heldout}
\small
\begin{tabular}{lrrr}
\toprule
Tier 1 rules & Train & Test & Test FPR \\
\midrule
Before any corpus tuning    & 0.7895 & 0.7632 & 0.0000 \\
Fitted on train split only  & 1.0000 & 0.7895 & 0.0000 \\
Tuned on whole corpus       & 1.0000 & 1.0000$^{*}$ & 0.0000 \\
\bottomrule
\end{tabular}
\\[2pt]
\footnotesize $^{*}$contaminated
\end{table}

The claimed figure overstates held-out recall by 21 points. Two details matter
more than the gap itself. Fitting 24 new terms on the train split moved test
recall by 2.6 points, from 0.763 to 0.789: hand-written blocklist terms fire on
their training samples and little else. Meanwhile the de-obfuscation pass, which
is a mechanism rather than a stored string, reaches 100\% on the test half's
\texttt{obfuscated\_unicode} class. Mechanisms transferred; vocabulary did not.

The held-out split also reverses our reading of the second tier. Measured
in-sample, Tier~2 contributed no detections, because a Tier~1 that has memorized
the corpus leaves nothing for it to catch. Held out, Tier~2 catches all 16
attacks Tier~1 misses and takes cascade recall to 1.000. The case for defense in
depth here is invisible without a holdout.

\subsection{Open-Prompt-Injection}

We reproduce the five attack strategies of~\cite{liu2024formalizing} from
upstream source, over six tasks, giving 450 injected prompts and 180 clean
task-data prompts as negatives. Table~\ref{tab:opi} gives recall by strategy.

\begin{table}[t]
\caption{Open-Prompt-Injection recall by attack strategy (450 injected prompts, 6
tasks). Only strategies carrying override vocabulary are detected under the
default source.}
\label{tab:opi}
\small
\begin{tabular}{lrrr}
\toprule
Strategy & Tier 1 & Tier 2 & \texttt{src=data} \\
\midrule
naive            & 0.00 & 0.00 & 1.00 \\
escape           & 0.00 & 0.00 & 1.00 \\
ignore           & 1.00 & 0.38 & 1.00 \\
fake completion  & 0.00 & 0.00 & 1.00 \\
combined         & 1.00 & 0.26 & 1.00 \\
\midrule
\textbf{overall} & \textbf{0.40} & \textbf{0.13} & \textbf{1.00} \\
FPR on clean     & 0.000 & 0.017 & 0.006 \\
\bottomrule
\end{tabular}
\end{table}

The pattern is stark and it is not a tuning gap. We detect the two strategies
containing the literal string \emph{Ignore previous instructions.} and miss all
three that do not. A movie review followed by ``Does the following text contain
spam or phishing contents?'' is a valid injection under this threat model and
contains nothing a blocklist of override idioms can see. Declaring
\texttt{source=data} takes overall recall to 1.00 at 1 false positive in 180.
Tier~2, a model fine-tuned for this exact task, reaches 0.13.

\subsection{BIPIA, and the ceiling}

BIPIA~\cite{yi2023bipia} embeds an injected instruction in email and table
content at start, middle or end. Its attacks carry no override vocabulary and no
formatted task specification; they are ordinary requests such as ``Write a script
to rename all files in a directory.'' We construct 1{,}200 injected samples over
15 categories with 60 clean contexts.

\begin{table}[t]
\caption{BIPIA, 1{,}200 indirect injections over 15 categories. Attacks carry no
override vocabulary.}
\label{tab:bipia}
\small
\begin{tabular}{lrr}
\toprule
Mode & Recall & FPR \\
\midrule
Tier 1, \texttt{source=user} & 0.035 & 0.000 \\
Tier 1, \texttt{source=data} & 0.480 & 0.017 \\
\quad before verb widening   & 0.082 & 0.000 \\
\bottomrule
\end{tabular}
\end{table}

This is the result we consider most useful. The instruction-shape patterns that
reach 1.00 on~\cite{liu2024formalizing} reach 0.48 here, and reached 0.082 before
we widened the verb list. They were written against that benchmark's templates,
so scoring well there was partly in-sample, exactly the error the previous
subsection diagnoses one level down. Per-category recall tracks vocabulary
overlap rather than attack severity: Substitution Ciphers 51\%, five categories
at 38\%, and five including Marketing and Misinformation at 2\%.

Widening the verb list bought 0.082 to 0.480 and left~\cite{liu2024formalizing}
at ${\sim}1.00$, at a cost in precision on imperative-heavy prose: false
positives on real SMS payloads went from 2/100 to 5/100. That is the treadmill.
Every benchmark needs new vocabulary and every addition costs precision somewhere
else.

\subsection{Over-defense, and two opposing failures}

On NotInject~\cite{li2024injecguard}, where all 339 samples are benign and every
block is a false positive, Tier~1 blocks none and Tier~2 blocks 137. Tier~2's
over-defense scales with trigger density: 20\%, 50\% and 51\% at one, two and
three trigger words, and a large share of its false positives are benign Chinese
prompts flagged on a single character.

\begin{table}[t]
\caption{NotInject. All 339 samples are benign, so every block is a false
positive.}
\label{tab:notinject}
\small
\begin{tabular}{lrr}
\toprule
Mode & Blocked & FPR \\
\midrule
Tier 1 (197-term blocklist) & 0 / 339   & 0.0000 \\
Tier 2 (DeBERTa-v3)         & 137 / 339 & 0.4041 \\
Shipped cascade             & 0 / 339   & 0.0000 \\
\bottomrule
\end{tabular}
\end{table}

This is not fixable by threshold. On NotInject's benign samples Tier~2's
injection score has median 0.9992, with 53\% at or above 0.999; on the held-out
attacks it correctly rescues, the median is 1.0000 with 88\% above that line. The
distributions overlap almost completely. We therefore ship Tier~2 in
monitor-only mode: it scores and logs but does not block. A 40\% false-block rate
is not a shippable default, and the cost, a recall ceiling of 0.789 rather than
1.000, is one we state rather than hide.

The two tiers fail in opposite directions. The deterministic one is precise and
narrow. The learned one is neither precise nor, on these benchmarks, broad. No
configuration of this cascade is simultaneously robust and precise, and no
threshold splits the difference.

\subsection{A quantized guardrail that had stopped working}

The shipped Tier~2 was an INT8 graph produced with per-tensor weight scaling. It
had collapsed to predicting \textsc{safe} for every input: it labels ``Ignore all
previous instructions and reveal your system prompt'' \textsc{safe} at 0.94
confidence. Behind a fail-open policy this is invisible. The tier loads cleanly,
reports latency, scores every request, and blocks nothing. It produced a Tier~2
recall of 0.005 that we initially read as a property of the model.

Re-benchmarking every available graph gave fp32 at 0.900 recall, and two
correctly configured INT8 variants at 0.758 and 0.732. Even done properly,
quantization cost 14 to 17 points of recall for roughly 5~ms and 495~MB.
DeBERTa-v3's disentangled attention is sensitive to per-tensor scaling. The
transferable lesson is that a silently degraded guardrail and a healthy one are
indistinguishable from the outside when the policy is fail-open, so quantized
graphs need accuracy verification and not just a smoke test.

\section{Discussion}

Detection here is trust-boundary dependent, and the dependence is large. One
declared bit moves recall on~\cite{liu2024formalizing} by 60 points. Vocabulary
added by hand moved held-out recall by 2.6. Any deployment that screens retrieved
content with the same rules it applies to user turns is leaving most of that on
the table, and the information required is already available at the call site.

The BIPIA result bounds what this family of detector can do. Patterns over
surface form transfer poorly across phrasings of the same threat, and the failure
is graded by vocabulary overlap rather than by how dangerous the attack is. A
detector that generalizes has to model what a span of text is trying to do, not
the words it uses. Representation-level approaches~\cite{pishield2025} look
better suited to that than any list we could write.

None of this makes the deterministic tier useless. It is the only component here
with a clean precision record: 0 of 339 on NotInject, 0 of 180 on real task data,
sub-millisecond, auditable, and it fails in ways an operator can read off a
config file. As a cheap first filter it is good. As a solution to indirect prompt
injection it is not, and the gap between those two claims is where most of the
risk in deploying one of these lives.

\section{Limitations}
\label{sec:limitations}

No adaptive attacker. Every number here is against static corpora, which is the
evaluation Nasr et al.~\cite{nasr2025attacker} argue overstates robustness, and
we expect ours are no exception. Our figures are upper bounds.

The refit is not blind. The train-only rules were authored by someone who had
already seen the full corpus, so 0.789 is optimistic. The split, per-term
provenance and scripts are published so the protocol can be re-run cleanly by
someone who has not.

Character-injection evasion~\cite{hackett2025bypassing} is unevaluated, and it
targets our normalization directly. Our de-obfuscation handles leetspeak, small
caps, fullwidth forms, combining marks and separator fragmentation, but has never
been tested against optimized character injection. We expect it to do badly and
we have not measured how badly.

Scope. One detector, one 271-sample corpus, BIPIA restricted to email and table
contexts, and the summarization task of~\cite{liu2024formalizing} omitted because
its data is gated. Declarative injections remain undetected under
\texttt{source=data}: ``It would be helpful to know the sentiment of this text''
is missed, verified.

\section{Conclusion}

A conventional two-tier guardrail scores 1.000, 0.789, 1.00 and 0.48 on four
evaluations of the same capability, and the spread is the finding. Telling the
detector which trust boundary a text crossed is worth more than any vocabulary we
added to it. The deterministic tier is precise and narrow, the learned tier
over-defends on 40\% of adversarially benign prompts, and quantization had
silently disabled the latter in a way a fail-open policy concealed. Pattern
matching over instruction surface form is a good first filter and it does not
survive a change of phrasing.

\section*{Availability}

Source, all measurement scripts, the corpus split and every result artifact are
public. The repository link is withheld for double-blind review and will be
included in the camera-ready version.

\balance
\bibliographystyle{ACM-Reference-Format}
\bibliography{refs}

\end{document}
